\vspace{78mm}
\centering
\section*{\underline{Skeletttheorie}}

\begin{tikzpicture}[remember picture, overlay]

% ---- ERSTE REIHE: Am Seitenursprung starten
\setlength{\rowheight}{30mm}

    % A1) Erste Kachel unten anbauen
    \setlength{\cellwidth}{98mm}
    \Tile{A1}{([xshift=10mm,yshift=-112mm]current page.north west)}{\cellwidth}{\rowheight}
        {
            \vspace{-\TileSep}
            \titlerm{Koeffizienten}\\ \vspace{-3mm}
            \begin{minipage}{0.42\cellwidth-\TileSep}
                \centering
                \begin{align*}
                    A_0 &= \alpha^* + \alpha_0 \\
                    A_1 &= 2\,\left(\alpha_2 - \alpha\right) = 4\,\frac{f}{t} \\
                    A_2 &= -3\,\alpha_2
                \end{align*}
            \end{minipage}%
            \begin{minipage}{0.58\cellwidth-\TileSep}
                \centering
                \begin{align*}
                    A_0 &= \alpha^* + \alpha_0 = \alpha^* - \frac{1}{\pi}\,\int\limits_0^\pi \frac{\text{d}z}{\text{d}x}\,\text{d}\Phi \\
                    A_n &= -\frac{2}{\pi}\,\int\limits_0^\pi\frac{\text{d}z}{\text{d}x}\,\cos\left(n\,\Phi\right)\,\text{d}\Phi
                \end{align*}
            \end{minipage}
        }{0}{1}{1}{0}

    % B1) Zweite Kachel rechts anbauen
    \setlength{\cellwidth}{40mm}
    \Tile{B1}{(A1.north east)}{\cellwidth}{\rowheight}
        {
            \vspace{-\TileSep}
            \titlerm{Kräfte}\\ \vspace{-3mm}
            \begin{align*}
                A &= G = \frac{\pi}{4}\,b\,\rho\,U_\infty\,A_1 \\[1.5ex]
                W_\text{i} &= \frac{\pi}{8}\,\rho\,\sum\limits_{n=1}^M\left(n\,A_n^2\right)
            \end{align*}
        }{0}{1}{1}{0}

    % C1) Dritte Kachel rechts anbauen
    \setlength{\cellwidth}{52mm}
    \Tile{C1}{(B1.north east)}{\cellwidth}{\rowheight}
        {
            \vspace{-\TileSep}
            \titlerm{Position der max. Wölbung}\\ \vspace{-2mm}
            \begin{equation*}
                \frac{\text{d}\bar{z}}{\text{d}\bar{x}} = 0~(\text{maximum von }\bar{z}(\bar{x}))
            \end{equation*}
            \begin{equation*}
                \frac{f}{t} = \frac{A_1}{4}
            \end{equation*}
        }{0}{0}{1}{0}

% ---- ZWEITE REIHE: Am Seitenursprung starten
\setlength{\rowheight}{38mm}

    % A2) Erste Kachel unten anbauen
    \setlength{\cellwidth}{95mm}
    \Tile{A2}{(A1.south west)}{\cellwidth}{\rowheight}
        {
            \titlerm{Beiwerte}\\ \vspace{3mm}
            \begin{minipage}{0.47\cellwidth-\TileSep}
                \centering
                \vspace{-3mm}
                \begin{align*}
                    c_a &= \pi\,\left(2\,A_0 + A_1\right)\\
                    \Delta c_p &= c_{p,\text{u}} - c_{p,\text{o}} = 2\,\frac{\gamma(\bar{x})}{U_\infty}
                \end{align*}
            \end{minipage}%
            \vline
            \begin{minipage}{0.53\cellwidth-\TileSep}
                \centering
                \vspace{-3mm}
                \begin{align*}
                    c_{m,\text{VK}} &= -\frac{\pi}{4}\,\left(2\,A_0 + 2\,A_1 + A_2\right) \\
                    c_{m,\sfrac{t}{4}} &= c_{m,\text{VK}} + \frac{c_a}{4}
                \end{align*}
            \end{minipage}
            \begin{equation*}
                \frac{\text{d}c_a}{\text{d}\alpha} = 2\,\pi~;\quad \frac{\text{d}c_m}{\text{d}\alpha} = 0
            \end{equation*}
        }{0}{1}{1}{0}

    % B2) Zweite Kachel rechts anbauen
    \setlength{\cellwidth}{95mm}
    \Tile{B2}{(A2.north east)}{\cellwidth}{\rowheight}
        {
            \titlebf{Nullauftrieb und -moment}\\ \vspace{2mm}
            \begin{minipage}{0.46\cellwidth-\TileSep}
                \centering
                \begin{align*}
                    \alpha_{A=0,\text{Skel}} &= -\frac{A_1}{2} - \frac{A_2}{3}\\
                    &= \alpha_{A=0,\text{Prof}} + \alpha_0 \\[1.0ex]
                    c_{a,\alpha=0} &= 4\,\pi\,\frac{f}{t}
                \end{align*}
                \vspace{-2mm}
            \end{minipage}%
            \vline
            \begin{minipage}{0.54\cellwidth-\TileSep}
                \centering
                \titlerm{Nullauftriebswinkel}\\ \vspace{-5mm}
                \begin{equation*}
                    c_a = 0\quad \Rightarrow A_0 = -\frac{A_1}{2}
                \end{equation*}
                \titlerm{Nullmoment}\\ \vspace{-5mm}
                \begin{equation*}
                    c_{,A=0} = -\pi\,\frac{f}{t} = -\frac{\pi}{4}\,\left(A_1 + A_2\right)
                \end{equation*}
            \end{minipage}
        }{0}{0}{1}{0}

\end{tikzpicture}